\renewcommand{\arraystretch}{1.4}

% =====================================================
% BLOCO LUA GLOBAL
% =====================================================

\begin{luacode*}

json = utilities.json
lfs  = require("lfs")

function contar_pendencias(codigo_tarefa)

  local pasta_issues = "tasks/issues"
  local total = 0

  for arquivo in lfs.dir(pasta_issues) do
    if arquivo:match("%.json$") and arquivo:match("^" .. codigo_tarefa) then

      local caminho = pasta_issues .. "/" .. arquivo
      local file = io.open(caminho, "r")

      if file then
        local conteudo = file:read("*a")
        file:close()

        local data = json.tolua(conteudo)

        if data.status == true then
          total = total + 1
        end
      end

    end
  end

  return total
end

function responsaveis()

  local pasta = "tasks"
  local responsaveis = {}

  -- ==========================
  -- LER TODOS OS JSON
  -- ==========================
  for arquivo in lfs.dir(pasta) do
    if arquivo:match("%.json$") then

      local caminho = pasta .. "/" .. arquivo
      local file = io.open(caminho, "r")

      if file then
        local conteudo = file:read("*a")
        file:close()

        local data = json.tolua(conteudo)

        if data.responsavel then

          local nome   = data.responsavel
          local codigo = arquivo:gsub("%.json$", "")
          local titulo = data.titulo or ""
          local peso   = tonumber(data.peso) or 0

          if not responsaveis[nome] then
            responsaveis[nome] = {
              tarefas = {},
              soma = 0
            }
          end

          local pendencias = contar_pendencias(codigo)

          table.insert(responsaveis[nome].tarefas, {
            codigo = codigo,
            titulo = titulo,
            peso   = peso,
            pendencias = pendencias
          })

          responsaveis[nome].soma =
            responsaveis[nome].soma + peso
        end
      end
    end
  end

  -- ==========================
  -- CRIAR LISTA ORDENADA
  -- ==========================

  local lista_nomes = {}

  for nome, _ in pairs(responsaveis) do
    table.insert(lista_nomes, nome)
  end

  table.sort(lista_nomes, function(a, b)

    if a == "Desativado" then
      return false
    elseif b == "Desativado" then
      return true
    else
      return a < b
    end

  end)

  -- ==========================
  -- IMPRESSÃO LATEX
  -- ==========================

  for _, nome in ipairs(lista_nomes) do

    local dados = responsaveis[nome]
    local soma_pendencias = 0

    -- Ordenar tarefas por código
    table.sort(dados.tarefas, function(a, b)
      return a.codigo < b.codigo
    end)

    tex.print("\\section*{" .. nome .. "}")

    tex.print("\\begin{longtable}{|p{0.15\\textwidth}|p{0.55\\textwidth}|p{0.1\\textwidth}|p{0.1\\textwidth}|}")
    tex.print("\\hline")
    tex.print("\\textbf{Código} & \\textbf{Descrição} & \\textbf{Peso} & \\textbf{Issues} \\\\ \\hline")
  
    for _, tarefa in ipairs(dados.tarefas) do
      soma_pendencias = soma_pendencias + tarefa.pendencias
      tex.print(
        tarefa.codigo .. " & " ..
        tarefa.titulo .. " & " ..
        tarefa.peso .. " & " ..
        tarefa.pendencias .. " \\\\ \\hline"
      )
    end

    tex.print("\\textbf{Soma} &  & \\textbf{" .. dados.soma .. "} & \\textbf{" .. soma_pendencias .. "} \\\\ \\hline")
    tex.print("\\end{longtable}")
    tex.print("\\bigskip")

  end

end


------------------------------------------------------
-- Função 1: Issues por prefixo (sem id no JSON)
------------------------------------------------------
function gerar_issues(prefixo)

  local resultados = {}

  for file in lfs.dir("tasks/issues") do
    if file:match("^" .. prefixo .. "%-(%d+)%.json$") then

      -- extrai número direto do nome do arquivo
      local numero = file:match("%-(%d+)%.json$")

      local f = io.open("tasks/issues/" .. file, "r")
      local content = f:read("*a")
      f:close()

      local data = json.tolua(content)

      if data.status == true then
        table.insert(resultados, {
          numero = tonumber(numero),
          descricao = data.descricao
        })
      end
    end
  end

  table.sort(resultados, function(a,b)
    return a.numero < b.numero
  end)

  tex.sprint("\\begin{tabular}{|p{0.1\\textwidth}|p{0.775\\textwidth}|}")
  tex.sprint("\\hline")
  tex.sprint("\\textbf{Issue} & \\textbf{Descrição} \\\\ \\hline")

  for _, item in ipairs(resultados) do
    tex.sprint(item.numero .. " & " .. item.descricao .. " \\\\ \\hline")
  end

  tex.sprint("\\end{tabular}")
  tex.sprint("\\bigskip")
  tex.sprint("\\\\")
end

------------------------------------------------------
-- Função 2: Ler tarefa pelo nome (sem id no JSON)
------------------------------------------------------
function read_task(nome)

  local caminho = "tasks/" .. nome .. ".json"
  local file = io.open(caminho, "r")
  local content = file:read("*a")
  file:close()

  local data = json.tolua(content)

  tex.sprint("\\newpage")
  tex.sprint("\\begin{tabular}{|p{0.1\\textwidth}|p{0.1\\textwidth}|p{0.57\\textwidth}|}")
  tex.sprint("\\hline")

  -- agora usamos o nome como ID
  tex.sprint(nome ..
    " & Peso: " .. data.peso ..
    " & " .. data.titulo ..
    " \\\\ \\hline")

  tex.sprint("\\multicolumn{2}{|p{0.2\\textwidth}|}{\\textbf{Responsável:}} & " ..
    data.responsavel ..
    " \\\\ \\hline")

  tex.sprint("\\multicolumn{2}{|p{0.2\\textwidth}|}{\\textbf{Estagiários(as):}} & " ..
    table.concat(data.estagiarios, "; ") ..
    " \\\\ \\hline")

  tex.sprint("\\multicolumn{3}{|p{0.8\\textwidth}|}{" ..
    data.descricao ..
    "} \\\\ \\hline")

  tex.sprint("\\end{tabular}")
  tex.sprint("\\bigskip")

  tex.sprint("\\\\ Pendências em aberto para " .. data.titulo .. ":\\\\ \\\\")

  gerar_issues(nome)
end

function list_json_by_prefix(prefix)
  local lfs = require("lfs")
  local files = {}

  for file in lfs.dir("tasks") do
    if string.match(file, "^"..prefix.."%d+%.json$") then
      table.insert(files, file)
    end
  end

  -- Ordenação numérica correta
  table.sort(files, function(a, b)
    local numA = tonumber(string.match(a, "%d+"))
    local numB = tonumber(string.match(b, "%d+"))
    return numA < numB
  end)

  return files
end

function loadjson(prefix)
  local files = list_json_by_prefix(prefix)

  for _, file in ipairs(files) do
    local id = string.gsub(file, "%.json$", "")
    tex.print("\\inputjson{" .. id .. "}")
  end
end

\end{luacode*}

% =====================================================
% COMANDOS LATEX
% =====================================================

\newcommand{\inputjson}[1]{\directlua{read_task("#1")}}
%\newcommand{\issues}[1]{\directlua{gerar_issues("#1")}}
\newcommand{\responsaveis}{%
  \directlua{responsaveis()}%
}
\newcommand{\loadjson}[1]{\directlua{loadjson("#1")}}